% !TEX TS-program = XeLaTeX
% use the following command:
% all document files must be coded in UTF-8
\documentclass[portuguese]{textolivre}
% build HTML with: make4ht -e build.lua -c textolivre.cfg -x -u article "fn-in,svg,pic-align"

\journalname{Texto Livre}
\thevolume{19}
%\thenumber{1} % old template
\theyear{2026}
\receiveddate{\DTMdisplaydate{2025}{9}{28}{-1}} % YYYY MM DD
\accepteddate{\DTMdisplaydate{2025}{10}{26}{-1}}
\publisheddate{\DTMdisplaydate{2026}{5}{2}{-1}}
\corrauthor{Lovania Roehrig Teixeira}
\articledoi{10.1590/1983-3652.2026.62011}
%\articleid{NNNN} % if the article ID is not the last 5 numbers of its DOI, provide it using \articleid{} commmand 
% list of available sesscions in the journal: articles, dossier, reports, essays, reviews, interviews, editorial
\articlesessionname{articles}
\runningauthor{Lima et al.} 
%\editorname{Leonardo Araújo} % old template
\sectioneditorname{Daniervelin Pereira~\orcid{0000-0003-1861-3609}}
\layouteditorname{Saula Cecília~\orcid{0009-0006-3069-8480}}

\title{O turno da Inteligência Artificial e o turno do professor na correção de textos da plataforma educacional ``Redação Paraná''}
\othertitle{The Artificial Intelligence shift and the teacher's shift in correcting texts on the ``Redação Paraná'' educational platform}
% if there is a third language title, add here:
%\othertitle{Artikelvorlage zur Einreichung beim Texto Livre Journal}

\author[1]{Agatha Suelin de Lima~\orcid{0009-0007-7132-5741}\thanks{Email: \href{mailto:agathalima@alunos.utfpr.edu.br}{agathalima@alunos.utfpr.edu.br}}}
\author[1]{Maisa Morelatto~\orcid{0009-0003-9537-6499}\thanks{Email: \href{mailto:maisamorelatto@outlook.com}{maisamorelatto@outlook.com}}}
\author[1]{Vanda Mari Trombetta~\orcid{0000-0002-0077-2794}\thanks{Email: \href{mailto:vandatrombetta@utfpr.edu.br}{vandatrombetta@utfpr.edu.br}}}
\author[1]{Lovania Roehrig Teixeira~\orcid{0000-0001-9614-8648}\thanks{Email: \href{mailto:lovaniateixeira@utfpr.edu.br}{lovaniateixeira@utfpr.edu.br}}}
\affil[1]{Universidade Tecnológica Federal do Paraná, Departamento de Letras, Pato Branco, PR, Brasil.}

\addbibresource{article.bib}
% use biber instead of bibtex
% $ biber article

% used to create dummy text for the template file
\definecolor{dark-gray}{gray}{0.35} % color used to display dummy texts
\usepackage{lipsum}
\SetLipsumParListSurrounders{\colorlet{oldcolor}{.}\color{dark-gray}}{\color{oldcolor}}

% used here only to provide the XeLaTeX and BibTeX logos
\usepackage{hologo}

% if you use multirows in a table, include the multirow package
\usepackage{multirow}

% provides sidewaysfigure environment
\usepackage{rotating}

% CUSTOM EPIGRAPH - BEGIN 
%%% https://tex.stackexchange.com/questions/193178/specific-epigraph-style
\usepackage{epigraph}
\renewcommand\textflush{flushright}
\makeatletter
\newlength\epitextskip
\pretocmd{\@epitext}{\em}{}{}
\apptocmd{\@epitext}{\em}{}{}
\patchcmd{\epigraph}{\@epitext{#1}\\}{\@epitext{#1}\\[\epitextskip]}{}{}
\makeatother
\setlength\epigraphrule{0pt}
\setlength\epitextskip{0.5ex}
\setlength\epigraphwidth{.7\textwidth}
% CUSTOM EPIGRAPH - END

% to use IPA symbols in unicode add
%\usepackage{fontspec}
%\newfontfamily\ipafont{CMU Serif}
%\newcommand{\ipa}[1]{{\ipafont #1}}
% and in the text you may use the \ipa{...} command passing the symbols in unicode

% LANGUAGE - BEGIN
% ARABIC
% for languages that use special fonts, you must provide the typeface that will be used
% \setotherlanguage{arabic}
% \newfontfamily\arabicfont[Script=Arabic]{Amiri}
% \newfontfamily\arabicfontsf[Script=Arabic]{Amiri}
% \newfontfamily\arabicfonttt[Script=Arabic]{Amiri}
%
% in the article, to add arabic text use: \textlang{arabic}{ ... }
%
% RUSSIAN
% for russian text we also need to define fonts with support for Cyrillic script
% \usepackage{fontspec}
% \setotherlanguage{russian}
% \newfontfamily\cyrillicfont{Times New Roman}
% \newfontfamily\cyrillicfontsf{Times New Roman}[Script=Cyrillic]
% \newfontfamily\cyrillicfonttt{Times New Roman}[Script=Cyrillic]
%
% in the text use \begin{russian} ... \end{russian}
% LANGUAGE - END

% EMOJIS - BEGIN
% to use emoticons in your manuscript
% https://stackoverflow.com/questions/190145/how-to-insert-emoticons-in-latex/57076064
% using font Symbola, which has full support
% the font may be downloaded at:
% https://dn-works.com/ufas/
% add to preamble:
% \newfontfamily\Symbola{Symbola}
% in the text use:
% {\Symbola }
% EMOJIS - END

% LABEL REFERENCE TO DESCRIPTIVE LIST - BEGIN
% reference itens in a descriptive list using their labels instead of numbers
% insert the code below in the preambule:
%\makeatletter
%\let\orgdescriptionlabel\descriptionlabel
%\renewcommand*{\descriptionlabel}[1]{%
%  \let\orglabel\label
%  \let\label\@gobble
%  \phantomsection
%  \edef\@currentlabel{#1\unskip}%
%  \let\label\orglabel
%  \orgdescriptionlabel{#1}%
%}
%\makeatother
%
% in your document, use as illustraded here:
%\begin{description}
%  \item[first\label{itm1}] this is only an example;
%  % ...  add more items
%\end{description}
% LABEL REFERENCE TO DESCRIPTIVE LIST - END


% add line numbers for submission
%\usepackage{lineno}
%\linenumbers

\begin{document}
\maketitle

\begin{polyabstract}
\begin{abstract}
 Este artigo analisa a plataforma ``Redação Paraná'', destacando a divisão das etapas de correção textual entre inteligência artificial (IA) e professores. O objetivo do estudo é compreender como o uso da plataforma impacta no processo de ensino-aprendizagem da produção textual, observando os tipos de correção disponíveis, sua adequação às concepções teóricas de revisão de textos e os impactos para alunos e docentes. Essa pesquisa se justifica pela relevância da produção textual como prática essencial no desenvolvimento de competências linguísticas e formação crítica dos estudantes e pelo aumento no uso de plataformas educacionais no contexto escolar. A correção, tarefa historicamente complexa e trabalhosa para o professor, ganha novas possibilidades com o uso da IA. Contudo, é necessário avaliar em que medida essa intervenção contribui para a reflexão sobre a escrita e para o aprimoramento das produções, evitando reduzir a atividade a um simples apontamento mecânico de erros. Os resultados indicam que a IA realiza predominantemente correções classificatórias e indicativas, oferecendo \textit{feedback} imediato. Apesar disso, apresenta limitações, como a incapacidade de detectar todos os erros ou explicar detalhadamente suas ocorrências. O professor, por sua vez, assume papel central ao corrigir aspectos discursivos e subjetivos, utilizando estratégias como comentários interativos e devolutivas, favorecendo a reescrita crítica e reflexiva do aluno. Conclui-se que a ``Redação Paraná'' contribui para otimizar o trabalho docente e apoiar o aprendizado, mas não substitui a mediação pedagógica, sendo imprescindível a atuação do professor para garantir avanços qualitativos na escrita.

\keywords{Redação Paraná\sep Correção\sep Inteligência Artificial\sep Professor}
\end{abstract}

\begin{english}
\begin{abstract}
This paper analyzes the ``Redação Paraná'' platform, highlighting the division of text correction stages between artificial intelligence (AI) and teachers. The objective is to understand how the use of the platform impacts the teaching-learning process of text production, observing the types of correction available, their adequacy to theoretical concepts of text revision, and the impacts on students and teachers. This research is justified by the relevance of text production as an essential practice in the development of students’ linguistic skills and critical thinking, and by the increased use of educational platforms in schools. Proofreading, a historically complex and laborious task for teachers, gains new possibilities with the use of AI. However, it is necessary to evaluate the extent to which this intervention contributes to reflection on writing and to the improvement of production, avoiding reducing the activity to a simple mechanical pointing out of errors. The results indicate that AI predominantly performs classificatory and indicative corrections, offering immediate feedback. Despite this, it has limitations, such as the inability to detect all errors or explain their occurrence in detail. The teacher, in turn, plays a central role in correcting discursive and subjective aspects, using strategies such as interactive comments and feedback, encouraging students’ critical and reflective rewriting. It can be concluded that the ``Redação Paraná'' contributes to optimizing teaching and supporting learning, but it does not replace pedagogical mediation. The teacher’s role is essential to ensure qualitative improvements in writing. 

\keywords{Redação Paraná\sep Correction\sep Artificial Intelligence\sep Teacher}
\end{abstract}
\end{english}
% if there is another abstract, insert it here using the same scheme
\end{polyabstract}

\section{Introdução}\label{sec-intro}
Desde 2020, devido à pandemia de Covid-19\footnote{Em 11 de março de 2020, a COVID-19 foi caracterizada pela OMS como uma pandemia, devido à ampla distribuição geográfica da doença no mundo. Disponível em: \url{https://www.paho.org/pt/historico-da-emergencia-internacional-covid-19}. Acesso em: 1 set. 2025.}, há um crescente uso de plataformas e número de aplicativos implantados pela Secretaria Estadual de Educação do estado do Paraná (SEED/PR), tais como ``Matemática Paraná'' (Matific), ``Inglês Paraná'', ``Leia Paraná'', ``Desafio Paraná'', entre outras \cite{morelatto2023}.

No que se refere ao componente curricular de Língua Portuguesa, que engloba práticas de linguagem como produção textual, leitura, oralidade e análise linguística \cite{parana2025a, parana2021b, parana2025b}, destaca-se a plataforma ``Redação Paraná'', objeto de discussão deste artigo. Esta plataforma tem o objetivo de apoiar o processo de ensino-aprendizagem da produção textual \cite{parana2023}. Um dos aspectos que chamam a atenção nessa plataforma é o fato da correção das produções dos alunos ser realizada, primeiramente, por intermédio da inteligência artificial (IA) e, somente na sequência, pelo professor.

A partir dessa observação, este artigo tem como objetivo explorar o funcionamento da plataforma ``Redação Paraná'' e refletir sobre as formas de correção utilizadas por ela, isto é, a correção no turno do professor e a correção no turno da IA. Uma discussão nesses moldes se justifica pela importância da realização de atividades de produção textual para o desenvolvimento das competências linguísticas dos estudantes e, assim, para a formação de cidadãos críticos e atuantes na sociedade. Ainda, compreender e avaliar as maneiras de correção das produções disponíveis na plataforma e até onde a atuação da IA é adequada faz-se necessário para que o processo de ensino de produção textual, que possui base teórica consolidada, seja realizado de maneira efetiva e adequada dentro do estável processo exposto pela literatura \cite{ruiz2015, conceicao2004}.

Nessa esteira, as reflexões realizadas terão como base os diferentes tipos de correções apontados em \textcite{ruiz2015} e \textcite{conceicao2004}. Partindo disso, é possível apontar pontos de contato e de afastamento em relação ao que a plataforma ``Redação Paraná'' propõe e, desse modo, refletir sobre como ela pode colaborar para o processo de ensino e aprendizagem da produção textual em sala de aula.

De modo geral, \textcite{ruiz2015} descreve quatro diferentes possibilidades de realizar uma correção textual no turno do professor: indicativa, classificatória, resolutiva e textual-interativa. Nossa hipótese é a de que a correção realizada pela IA se enquadre em pelo menos uma dessas quatro. Ainda, por ser uma tecnologia em uso nas escolas de todo o sistema estadual de ensino do Paraná, espera-se que a ``Redação Paraná'' proporcione possibilidades de trabalho com o texto do aluno que valorizem a construção do conhecimento sobre a escrita, possibilite uma revisão que favoreça a reflexão e a aprendizagem dos estudantes e também colabore com o trabalho de revisão do professor. Nesse contexto, supõe-se que as atividades e as correções da plataforma sejam claras e levem os estudantes a refletirem sobre sua escrita, bem como possibilitem a reescritura das produções textuais de forma crítica e ativa, para que, assim, a ``Redação Paraná'' não seja somente um instrumento de resolução de problemas gramaticais, especialmente no turno da IA.

A fim de discutir esses elementos, este artigo, primeiramente apresenta as bases teóricas relacionadas ao processo de correção ou revisão de textos; depois, apresenta alguns aspectos práticos de uso da plataforma nesse processo, ao mesmo tempo que realiza uma análise crítica, comparando o que a literatura defende e a prática de correção possível na plataforma; na sequência, são expostas as considerações finais.

\section{Processo de correção ou revisão de textos: o que diz a literatura?}
Tendo em vista que este artigo consiste em uma reflexão sobre os tipos de correção disponibilizados na plataforma ``Redação Paraná'' e seus impactos no ensino-aprendizagem, nesta seção, apresentam-se os conceitos de \textcite{ruiz2015} e \textcite{conceicao2004} sobre os tipos de correção de produção textual. 

Apesar da importância da produção textual como atividade que leva ao desenvolvimento das habilidades linguísticas dos alunos, ela também é uma das atividades que mais onera o professor de Língua Portuguesa. A correção, como define \textcite[p. 19]{ruiz2015}, ``é o trabalho que o professor (visando a reescrita do texto do aluno) faz nesse mesmo texto, no sentido de chamar sua atenção para algum problema de produção''. Assim, é uma atividade que consiste em leituras atentas de cada texto em suas diferentes versões e uma análise da escrita de cada uma delas a fim de identificar inadequações em cada um dos contextos, além de uma análise holística do progresso dessa escrita e da evolução de cada aluno. Se se considerar uma turma de 30 alunos, o professor lerá no final de um único processo completo de produção de um texto, envolvendo somente uma escrita e uma reescrita, um total de 60 versões de texto. Se se multiplicar esse valor por, no mínimo, um texto por semana, em um mês o professor terá lido 240 textos em uma única turma.

Ainda, neste artigo e assim no processo de correção realizado pelo professor, corrobora-se a concepção de \textcite[p. 96]{goncalves2023}, de que o texto ou a produção textual deve ser considerada numa perspectiva interacionista. Nesse sentido, um texto produzido na escola também possui conteúdos que dialogam com outros textos, pois apresenta a posição discursiva do enunciador-aluno que revive vozes de outros enunciados que estão relacionados ao contexto sócio-histórico e ideológico de escrita. Afinal, ``[...] Todas as palavras (enunciados, produções de discurso e literárias), além das minhas próprias, são palavras do outro. Eu vivo em um mundo de palavras do outro'' \cite[p. 379]{bakhtin2011}. Ainda, pelo viés dialógico da linguagem, o texto produzido pelo aluno é uma palavra que espera uma contrapalavra \cite{bakhtin2002, bakhtin2010, bakhtin2011} de seu leitor, o professor. Nesse sentido, o escritor/enunciador aguarda um acontecimento responsivo do leitor/enunciatário que receberá o que se chama de ``palavras alheias''. Segundo \textcite[p. 23]{ruiz2015}, nessa esteira da concepção da contrapalavra, o ``ouvinte/leitor não é absolutamente um `receptor' passivo, já que lhe cabe atuar sociocognitivamente sobre o material linguístico de que dispõe''.

Cabe ressaltar, então, que a contrapalavra do professor vem, muitas vezes, a partir do lugar social que ocupa nesse processo, isto é, da posição de leitor extremamente crítico e poucas vezes colaborativo. Essa contrapalavra, assim, se diferencia da contrapalavra de um leitor comum, pois vem acompanhada pela responsabilidade de corrigir o texto e pela necessidade de aprimorar as habilidades linguísticas dos alunos, os autores desses textos.

Mas não se pode esquecer que antes da leitura do professor e da esperada contrapalavra, o texto em produção recebe uma leitura do seu próprio produtor em que ele monitora sua escrita e reflete sobre ela. Nesse momento, ``[...] quem aponta os problemas do texto não é o produtor, mas o leitor que convive no produtor, ou o produtor desempenhando o papel de leitor, ainda que se trate do mesmo indivíduo; é, pois, o leitor-produtor'' \cite[p. 20]{ruiz2015}.

Quando se menciona a tarefa de correção do texto realizada pelo professor, \textcite[p. 34]{ruiz2015} indica ainda que ``[...] o texto interventivo do professor é um texto sobreposto ao texto do aluno, isto é, um texto que se produz `na carona' deste''. Assim sendo, o texto interventivo do professor pode estar em três regiões do espaço compartilhado com o aluno: (i) no corpo; (ii) margem; ou (iii) em sequência ao texto do aluno (``pós-texto'') \cite[p. 35]{ruiz2015}. No que se refere à correção, \textcite{ruiz2015} elenca quatro estratégias de correção textual que podem ser utilizadas pelos professores: indicativa, resolutiva, classificatória e textual-interativa. É importante notar que as formas de correção não são excludentes, segundo \textcite[p. 35]{ruiz2015}, e ocorrem de um modo ``híbrido'', pois ``[...] os professores, na maior parte dos casos, utilizam-se de mais de uma forma interventiva, mesclando tipos diferentes de correção''.

Na correção indicativa, o professor indica as palavras, as expressões ou as frases com problemas por meio de alguma sinalização, verbal ou não, na margem do texto do aluno ou então no corpo da redação \textcite[p. 36]{ruiz2015}. Nesse sentido, os apontamentos são realizados somente para indicar os problemas, que se resumem às irregularidades lexicais ou ortográficas, e não há esclarecimentos sobre o quê e o porquê de eles terem de ser corrigidos pelos alunos.

A correção resolutiva ``[...] é uma tentativa do professor assumir, pelo aluno, a reformulação de seu texto'' \cite[p. 41]{ruiz2015} e há diferentes maneiras de apresentar ao aluno uma solução para o problema encontrado, seja acrescentando, retirando, substituindo ou mudando de lugar partes do texto. Segundo \textcite{ruiz2015}, a grande concentração de resolutivas se dá no corpo do texto, sendo pouco frequentes as resoluções na margem e menos ainda no ``pós-texto''. Como se observa, essa forma de correção requer mais tempo de análise e maior esforço do professor, já que ele realiza a correção pelo aluno. Assim, o aluno é passivo nesse processo, pois não há chance de ele realizar um trabalho de reflexão, já que a adequação já vem indicada.

A correção classificatória consiste na identificação não ambígua dos erros por meio do uso de símbolos escolhidos pelo professor para que o aluno corrija seu erro de maneira autônoma. \textcite[p. 45]{ruiz2015} indica que há utilização de um certo conjunto de símbolos (letras ou abreviações), ``[...] escritos em geral à margem do texto, para classificar o tipo de problema encontrado. Essas letras, conhecidas dos alunos, fazem parte de um código de correção que varia de professor para professor''. Assim, a correção classificatória indica mais precisamente a natureza do problema encontrado no texto do aluno e permite que ele realize a correção por si, o que permite alguma reflexão por parte do aluno.

Na correção textual-interativa, o professor deixa comentários mais longos (``bilhetes'', conforme \textcite[p. 47]{ruiz2015}) que são geralmente após o texto do aluno, i.e., no ``pós-texto''. Um ponto relevante apontado pela autora neste tipo de correção é que ``esse distanciamento, visível no papel em termos espaciais, na verdade, reflete a troca de turnos que ocorre na interlocução aluno-produtor/professor-corretor/aluno-revisor'' \cite[p. 48]{ruiz2015}. Segundo a autora, esses comentários têm duas funções: (i) abordar a tarefa de revisão pelo aluno ou os problemas do texto; (ii) expor, metadiscursivamente, acerca da própria tarefa de correção pelo professor. Para fins de organização, as estratégias de correção de \textcite{ruiz2015} são sumarizadas na Tabela \ref{fig-1}.

%--- codigo da tabela figura 1 ---%
\begin{table}[h!]
\centering
\begin{threeparttable}
\caption{Estratégias de correção de acordo com \textcite{ruiz2015}.}\label{tab-1}
\begin{tabular}{p{3.5cm} p{10cm}}
\toprule
Estratégias de correção & Modo de correção \\
\midrule

Indicativa 
& - É o mero apontar dos erros, indicando-os no corpo do texto por meio de setas, círculos, riscos etc. \newline
  - Esta estratégia é voltada aos erros gramaticais ou lexicais. \\

\addlinespace

Resolutiva 
& - Consiste na reescrita, pelo professor, de palavras, frases ou períodos inteiros, por meio de estratégias de adição, substituição, deslocamento e de supressão. \\

\addlinespace

Classificatória 
& - Os professores criam um código (letras, abreviações ou sinais), conhecido pelos alunos, para indicar os erros. \newline
  - Essa estratégia faz com que o aluno busque corrigir seu próprio erro. \\

\addlinespace

Textual-interativa 
& - São comentários longos como ``bilhetes''. Nesta estratégia, o professor não está preocupado em discutir os erros do aluno, mas na tarefa de revisão. \\

\bottomrule
\end{tabular}
\source{Adaptação de \textcite{ruiz2015}.}
\end{threeparttable}
\end{table}

Ainda, sobre a tarefa de correção, \textcite{ruiz2015, conceicao2004} também afirmam que esta deve ir além da correção gramatical e incluir os problemas discursivos, isto é, os aspectos pragmáticos e semântico-conceituais do texto. Além disso, é importante ter em mente que o ato de corrigir um texto influencia de forma direta as próximas produções textuais dos estudantes, a forma como eles irão fazer a reescrita do texto e compreender os erros cometidos. Nesse sentido, a atitude adotada pelo professor na realização da atividade de correção textual é crucial:

\begin{quote}
    [...] o docente não deve apenas apontar os erros gramaticais, mas também se posicionar como leitor crítico e avaliar a escrita do discente, abrangendo os aspectos fundamentais ligados aos gêneros discursivos, iniciando assim, um processo interativo entre ambos \cite[p. 97]{goncalves2023}.
\end{quote}

Ainda afirmam as autoras \cite{ruiz2015, conceicao2004} que a correção do professor deve ser dialógica com o texto e com o autor, isto é, o professor deve ``[...] deslocar-se da posição de avaliador, de juiz do texto do aluno para tornar-se um interlocutor de verdade'' \cite[p. 328]{conceicao2004}. O professor como interlocutor questiona, elogia e incentiva o aluno a refletir sobre seu discurso e a reescrevê-lo a partir de seus apontamentos na correção. Essa forma de correção está relacionada com o entendimento da língua em sua vivacidade e dinamicidade, a língua como fator essencial para as interações sociais. A interlocução esperada no processo de produção textual, se baseia na ideia de que a língua tem essencialmente um caráter interativo, pois

\begin{quote}
    a língua não está apenas dicionarizada, mas sim viva nas interações sociais e historicamente situadas. Logo, a noção de língua viva é pertinente para que possamos estabelecer uma metodologia de correção de textos/enunciados, pois, no processo interativo, as relações de sentidos serão constituintes da produção do aluno \cite[p. 96]{goncalves2023}.
\end{quote}

Como afirma \textcite{ruiz2015}, a correção do professor visa à reescrita do texto, feita pelo aluno. A reescrita após os apontamentos é ``realizada pelo aluno em função de intervenções escritas do professor, via correção, com vistas a uma melhor legibilidade do seu texto'' \cite[p. 25]{ruiz2015}. Assim, ela é importante porque permite ao aluno uma evolução qualitativa do texto, visando a qualidade discursiva \cite{conceicao2004}.

Segundo \textcite[p. 215]{fuzer2012bilhete}, uma avaliação textual deve possibilitar que o estudante avance ``no conhecimento e na utilização dos recursos da língua materna em uso nos mais diversos contextos''. Nesse contexto, é na reescrita que esses avanços devem aparecer, por isso, o trabalho do professor deve ser analítico e holístico. Quando se comparam versões de texto e se observa uma progressiva capacidade de uso da língua, então, se obteve sucesso no processo de ensino-aprendizagem de produção textual.

Na seção seguinte, apresentam-se aspectos práticos da plataforma ``Redação Paraná'' relacionados à tarefa de correção no turno da IA e no turno do professor.

\section{Aspectos práticos de uso da plataforma ``Redação Paraná''}
A plataforma ``Redação Paraná'' é uma plataforma educacional implantada pela SEED/PR e disponível a alunos e a professores do Ensino Fundamental, anos finais, e Ensino Médio de todas as escolas públicas estaduais do Paraná. Esta plataforma tem como objetivo proporcionar aos estudantes a elaboração de textos em ambiente digital \cite{parana2023}, isto é, possibilita o desenvolvimento de redações online ``de forma integrada com o professor'' \cite{parana2023}. Nesse contexto, é um ambiente que reúne ensino-aprendizagem, já que nele o estudante produz um gênero textual e o professor realiza a correção desse texto, permitindo que o texto seja reescrito pelo aluno, se necessário, seguindo as indicações do professor.

De acordo com a SEED \cite{parana2023}, no entanto, há dois atores na correção realizada na plataforma ``Redação Paraná'': a IA e o professor. A primeira correção é feita pela IA, que é responsável por corrigir a ``estrutura da língua -- Gramática'' \cite{parana2023}. Após a verificação gramatical feita pela IA, o texto é disponibilizado ao professor que fica responsável por corrigir ``a parte discursiva e subjetiva da redação elaborada pelo aluno'' \cite{parana2023}.

Neste artigo, entendem-se essas duas maneiras de corrigir textos a partir do conceito de turnos, aspecto que se alia à concepção de correção textual de \textcite{ruiz2015}. Para a autora, a atividade de correção de um texto é feita em turnos que se revezam entre o professor, em uma primeira leitura, e o aluno, entre leituras e revisões posteriores deste texto. Uma particularidade da ``Redação Paraná'' é o acréscimo de um terceiro turno que é realizado pela IA, cuja correção é a primeira a ser realizada, antes até daquela realizada pelo professor. Deste modo, compreende-se que a plataforma possibilita dois turnos de correção, desconsiderando o do aluno: a IA realiza uma correção de viés gramatical e, depois, o professor analisa a discursividade e a subjetividade da redação \cite{parana2023}.

A IA (Gemini da empresa Google \cite{parana2025a}) utilizada para a correção na plataforma faz uso de mais de 2,5 mil regras de linguagem \cite{parana2021a}. Ela realiza, assim, uma correção automática, apontando os erros gramaticais cometidos no texto, isto é, indica problemas da superfície textual, deixando inadequações profundas, como questões discursivas, para o professor analisar. Nesse contexto, observa-se que a correção textual do professor, na plataforma, também vai além do apontamento de erros gramaticais nos textos dos alunos, o que vai ao encontro do que se afirma na literatura. A seguir, detalham-se as atividades de correção possíveis nos turnos da IA e do professor e apontam-se os aspectos positivos e limitantes em cada um deles.


\subsection{A correção no turno da IA}
A correção feita pela IA da plataforma ``Redação Paraná'' trata-se de uma correção prévia e os alunos podem visualizá-la durante o desenvolvimento do texto. Para isso, o aluno deve selecionar a opção de ``correção \textit{online}''. A partir disso, a IA localiza e indica os erros presentes no texto, destacando palavras ou sentenças que devem ser reescritas. A partir desses apontamentos, o aluno pode realizar as primeiras adequações em seu texto.

A correção \textit{online} feita pela IA se dá com base em inadequações de diferentes ordens que são classificadas e identificadas por cores diferentes: acentuação (amarelo), ortografia (verde), semântica (rosa), sintaxe (laranja) e ``outras'' (azul). Na Figura \ref{fig-1} pode ser observado que os problemas de diferentes naturezas são apontados por cores diferentes:

%--- codigo da figura imagem 1 ---%
\begin{figure}[h!]
\centering
\begin{minipage}{.95\textwidth}
\includegraphics[width=\textwidth]{Imagens/image1.png}
\caption{Correção \textit{online} da plataforma ``Redação Paraná''.}
\label{fig-1}
\source{Print de tela \cite{parana2021b}.}
\end{minipage}
\end{figure}

A correção realizada nessa etapa é semelhante ao que \textcite{ruiz2015} denomina ``correção classificatória''. De acordo com a autora, este tipo de correção é feito usando códigos de classificação desenvolvidos pelo professor e serve para identificar o tipo de erro cometido no texto pelo aluno. Nesse contexto, o professor sublinha uma palavra ou sentença e indica, por exemplo, a letra G, que poderia ser o código determinado para ``gramática''; assim, o aluno entende que naquele item há um erro gramatical.

A correção classificatória baseada na classificação dos erros existe na plataforma ``Redação Paraná'', a diferença é que ela não é feita por códigos, mas pela distinção de cores, como visto na Figura \ref{fig-1}, que corresponde aos elementos linguísticos analisados pela IA. Esse tipo de correção apresenta algumas vantagens, uma delas é:

\begin{quote}
    [...] esta é uma correção que respeita os princípios de uma boa correção porque prevalece uma atitude operativa (identifica-se o erro, especificando a operação que o aluno errou), enquanto nas outras duas\footnote{As outras duas correções mencionadas pela autora são a resolutiva e a indicativa.} prevalece uma atitude prescritiva (indica-se ou resolve-se o erro) \cite[p. 325]{conceicao2004}.
\end{quote}

Assim, é positivo o fato de a plataforma disponibilizar imediato \textit{feedback} aos alunos, não corrigindo seus erros, mas apenas indicando quais são, onde estão e sugerindo alterações. Entretanto, a correção classificatória proposta pela IA não apresenta clareza suficiente para que o aluno entenda os seus erros e tenha uma atitude ativa e operativa sobre eles. Nesse sentido, é importante mencionar que o professor também possui uma função durante o turno da IA, pois ela ``não é capaz de compreender o motivo do desvio ou mesmo se se trata efetivamente de um, sem falar das questões de coerência, projeto textual, qualidade da argumentação etc'' \cite[p. 48]{lima2023}.

Assim, a correção classificatória realizada pela plataforma é restrita e exige que o professor realize o seu turno de correção partindo de critérios gramaticais complementares (não contemplados pela IA) e aspectos linguísticos amplos, e ainda, explicite as dúvidas dos alunos, aspecto que a IA não consegue realizar.

Observa-se uma limitação do turno de correção da IA, a partir da análise da Figura \ref{fig-2}, abaixo: há algumas inadequações, que mesmo estando no rol de problemas abarcados pela IA, não foram apontados, tais como em ``uma etapas'' (circulado em vermelho por estas autoras) em que há um problema de concordância entre o artigo ``uma'', no singular, e o substantivo ``etapas'', no plural.

%---codigo da figura imagem 2 ---%
\begin{figure}[h!]
\centering
\begin{minipage}{.95\textwidth}
\includegraphics[width=\textwidth]{Imagens/Fig2.png}
\caption{Problema gramatical não rastreado pela IA.}
\label{fig-2}
\source{Print de tela \cite{parana2021b}.}
\end{minipage}
\end{figure}

Nesse primeiro momento da escritura do texto pelo aluno, antes do envio da sua produção textual, como visto, ele tem a primeira oportunidade de reescrever esse texto corrigindo os erros apontados pela IA. A tarefa de reescrita deve acontecer diversas vezes durante a produção textual, ``como [n]a pré-escritura, [n]a escritura e [na] pós-escritura'' \cite[p. 25]{ruiz2015}, pois corresponde aos momentos de reflexão sobre os instrumentos que a linguagem possui e os diferentes sentidos que podem ser veiculados com cada um deles. Até o envio do texto, a plataforma possibilita a pré-escritura e a escritura do texto pelo estudante em sua área de aluno.

É importante pontuar que se o aluno realiza a produção do texto na tela do computador e já nessa etapa a IA faz suas indicações, há uma perturbação da leitura que o próprio autor faz durante a produção de seu texto. O que ocorre é que o aluno não terá oportunidade de ler o texto sem o apoio da IA e assim não refletirá autonomamente sobre os instrumentos linguísticos disponíveis para o aprimoramento de sua escrita.

Como se pode ver na Figura \ref{fig-2}, no turno de correção da IA, abre-se uma caixa de diálogo quando o aluno clica sobre o erro apontado pela IA. Nessa caixa de diálogo há: uma justificativa para o apontamento de erro feito; uma explicação para aquela correção; e ainda, o termo corrigido. Aqui, um ponto relevante para o processo de aprendizagem é a explicação ofertada pelo sistema por meio dessa caixa de diálogo, no entanto, não é garantido que o estudante clique, leia e procure entender essa explicação. Nesse momento, há dois cenários possíveis: de um lado, um estudante que possui autonomia e interesse que irá ler e tentar compreender a explicação do porquê de sua escrita estar errada; e, do outro, um estudante que automaticamente irá clicar sobre o termo corrigido sem nenhum tipo de reflexão. Para o primeiro, o turno de correção da IA será positivo e representará uma oportunidade de aprendizado; já para o segundo, somente incentivará uma dependência do corretor automático dessa tecnologia.

Nessa etapa, um aspecto limitante é a nomenclatura de indicação dos problemas feita pela IA, que não fornece pistas mais detalhadas da inadequação encontrada no texto redigido pelo aluno; por exemplo, a indicação de problema com ``semântica'' é genérica e não é autoexplicativa. O ponto positivo, por sua vez, é o fato de o aluno ter explicações disponíveis para seu erro quando posicionar o cursor sobre a expressão indicada pela IA e selecionar a opção ``ver explicação'', mas, para isso, é preciso que o aluno tenha interesse em fazê-lo. Nesse ponto, observa-se uma correção indicativa sendo realizada pela IA, apontando os erros do aluno, mas também ocorre uma correção resolutiva, já que a IA pode fornecer a expressão corrigida ao aluno e também faz uma correção classificatória, categorizando os problemas com base em cinco tipos (acentuação, ortografia, semântica, sintaxe e ``outras'').

Após o envio da redação pelo aluno via ``Redação Paraná'', o professor recebe o texto corrigido pela IA também via plataforma. Se o aluno não realizou as correções indicadas pela IA no momento da produção, as inadequações remanescentes são apresentadas também ao professor (classificadas também como ortografia, semântica, sintaxe e outras). É importante notar que, se o aluno realizar a correção dos problemas apontados pela IA, não é enviado um relatório ao professor contendo o número e a natureza dos erros encontrados, o que se constitui em um aspecto limitante do turno de correção da IA na ``Redação Paraná''. A geração e a disponibilização do relatório de ``erros corrigidos no turno da IA'' ao professor em seu turno de correção possibilitariam que ele ficasse sabendo quais são os problemas gramaticais recorrentes encontrados durante a escrita dos textos de seus estudantes. Nesse sentido, permitiria que o docente diagnosticasse as dificuldades de práticas de linguagem e, assim, sanasse tais dificuldades com atividades de revisão ou de retomada de aspectos já trabalhados.

A partir dessa pré-correção realizada pela IA, a plataforma oferece uma nota de 0 a 40 pontos, que será somada à nota que o professor irá atribuir na sequência. No entanto, ressalta-se que a nota atribuída pela IA pode ser editada pelo professor \cite{parana2025a, parana2021b, parana2025b}, o que é um aspecto positivo do processo de correção da plataforma. Na sequência, apresentam-se aspectos da correção na plataforma ``Redação Paraná'' no turno do professor.

\subsection{A correção no turno do professor}
A leitura, a análise e a correção de textos são as tarefas mais onerosas do professor de Língua Portuguesa. Além disso, o texto é um produto (inacabado) de um processo longo, que começa com leituras de textos do mesmo gênero e segue com análises de suas características comuns. Só depois disso, inicia-se um planejamento de escrita pelo seu autor.

No contexto das instituições públicas de educação do Paraná, a plataforma ``Redação Paraná'' oferece possibilidades de realização dessas tarefas \textit{online}, mas este artigo não se aterá a esses pontos, tendo em vista que o seu escopo se dá sobre as etapas de correção na plataforma. Nesse contexto, uma vez que a IA, por seu turno, indicou as questões gramaticais (acentuação, ortografia, semântica, sintaxe e ``outras''), o professor volta seu foco a aspectos textuais mais complexos e profundos. De acordo com a SEED/PR \cite{parana2023}, fica como responsabilidade do professor a correção da ``estrutura do texto e a lógica entre as partes''.

O papel do professor, segundo \textcite[p. 218]{fuzer2012bilhete}, na correção textual é de auxiliar ``o escritor a usar ou expandir seu conhecimento sobre formas de escrever, focalizando as convenções do gênero discursivo, o uso adequado do sistema linguístico e, se necessário, a adequação do texto ao assunto proposto''. Nesse sentido, a plataforma \textcite{parana2023} propõe os seguintes critérios para guiar os elementos textuais que devem ser o foco durante o turno de correção do professor. Sendo eles\footnote{As rubricas para cada gênero textual podem ser consultadas em \url{https://professor.escoladigital.pr.gov.br/redacao_parana}.}:

\begin{enumerate}[label=\alph*)] 
\item adequação ao gênero, respeitando os aspectos discursivos do gênero proposto;
\item adequação ao gênero, a respeito da adequação linguística do gênero;
\item convenções de escrita;
\item marcas de autoria e tema.
\end{enumerate}

Como se vê, os critérios indicados pela plataforma sugerem que o professor esteja atento ao conteúdo e à estrutura do gênero textual que está sendo produzido pelo aluno. Ainda que mantenha também o foco nas convenções da escrita, aspectos que a IA em seu turno já deveria ter apontado e corrigido. Também que observe a subjetividade do autor no texto e sua contribuição ao tema.

Assim, partindo da correção prévia realizada pela própria plataforma, o professor deverá sanar as lacunas que podem ter sido deixadas pela IA. Nem sempre a IA dará conta de corrigir os aspectos a que se propõe ou de deixar explícitos os equívocos cometidos pelo aluno em seu texto \cite{lima2023}. Nesse sentido, o papel do professor é complementar essa correção, seguindo os critérios propostos na própria plataforma. Deste modo, a correção realizada pelo professor vai além do mero apontar de erros, que seria uma incumbência da IA, estendendo-se para a correção dos aspectos discursivos.

A correção classificatória feita pela IA é importante, mas é superficial, como afirma \textcite[p. 325]{conceicao2004}: ``embora a correção classificatória respeite os princípios de uma boa correção [...] se esta ficar presa à classificação dos problemas gramaticais, não tocará, nem de longe, no real problema do texto do aluno''. Desta forma, fica para o professor a complexa atividade de corrigir os problemas discursivos do texto, estar atento à subjetividade do autor nesse processo e revisar os aspectos gramaticais que permaneceram, mesmo após a correção da IA.

No turno do professor, observa-se que a plataforma tenta guiar a correção por meio dos quatro critérios acima citados, mas ela também apresenta perguntas-guia para cada um dos critérios indicados. Por exemplo, no critério ``a) adequação ao gênero, respeitando os aspectos discursivos do gênero proposto;'' há a seguinte pergunta: \emph{o texto apresenta os estágios composicionais típicos do conto: introdução; desenvolvimento; clímax e desfecho?} Assim, em cada um dos critérios de correção a plataforma questiona e leva o professor a refletir sobre as questões de adequação ao gênero textual, recursos linguísticos, coesão, coerência e organização das ideias.

Uma funcionalidade da plataforma, que se encaixa nos tipos de correção indicativa e resolutiva, durante o turno do professor, é a possibilidade de assinalar algum problema no texto do aluno, destacando uma palavra, expressão ou trecho do texto do aluno, e comentar esse problema ou fornecer uma correção por meio de um ``balão''. Isso permite uma correção indicativa, se o docente somente realizar a indicação do trecho que estiver com problemas, mas também permite uma correção resolutiva se o professor marcar a expressão e ainda colocar a forma corrigida no ``balão'' de comentário.

Ainda no turno do professor, ao final da tela, a plataforma disponibiliza uma caixa de texto destinada à escrita de considerações pelo professor, no campo denominado ``Devolutiva'', como se observa na Figura \ref{fig-3}. Nessa caixa de texto podem ser feitos pequenos comentários antes do envio ao aluno de seu texto corrigido. De acordo com \textcite{ruiz2015}, pode-se classificar uma correção nesses moldes como ``textual-interativa'', pois é uma

\begin{quote}
    [...] forma alternativa encontrada pelo professor para dar conta de apontar, classificar ou até mesmo resolver aqueles problemas da redação do aluno que, por alguma razão, ele percebe que não basta via corpo, margem ou símbolo \cite[p. 52]{ruiz2015}.
\end{quote}

%--- codigo da figura imagem 3 ---%
\begin{figure}[h!]
\centering
\begin{minipage}{.62\textwidth}
\includegraphics[width=\textwidth]{Imagens/Fig3.png}
\caption{Campo para envio de comentários do professor ao aluno.}
\label{fig-3}
\source{Print de tela \cite{parana2021b}.}
\end{minipage}
\end{figure}

Nessa estratégia de correção, a finalidade é estabelecer um diálogo entre o texto escrito pelo estudante e a tarefa de correção realizada pelo professor, esclarecer pontos sobre a revisão ou ainda propor tarefas sobre uma reescrita e correção de seu texto. O professor possui papel de interlocutor, de modo que dialoga com o texto e com o autor (o estudante), exercitando a habilidade de interação desse sujeito-leitor, aspecto necessário na atividade docente, pois nesse local pode-se dar conselhos e também ressaltar os pontos positivos que o aluno apresentou durante a escrita/reescrita do texto. A diferença é que, no modo \textit{online}, disponível na ``Redação Paraná'', isto não é feito por bilhetes, mas como comentário logo após o texto.


\textcite{ruiz2015} afirma que a correção textual-interativa tem como um de seus objetivos principais a reescrita do texto. Como visto anteriormente, a plataforma possibilita a reescrita do texto ainda durante a sua produção (pré-escritura e escritura), antes do envio ao professor, entretanto, se solicitada pelo professor, a pós-escritura pode acontecer na plataforma após a correção realizada pelo professor. O professor pode solicitar a reescrita quantas vezes achar necessário, o que \textcite{conceicao2004} aponta como ponto positivo na construção textual.

A partir da Figura \ref{fig-3}, também se observa que a nota atribuída pelo professor vai de 0 a 60 pontos, isto é, tem um peso maior do que aquele atribuído à correção feita pela IA. Esse aspecto é positivo, tendo em vista que no turno da IA ocorre uma correção superficial, pois seu foco é gramatical, enquanto no turno do professor a correção se amplifica para os sentidos e os aspectos discursivos do gênero textual, além do progresso apresentado pelo aluno, se estivermos considerando uma reescrita.

\section{Considerações Finais}
A correção textual é uma das atividades mais desafiadoras do trabalho do professor de língua portuguesa \cite{conceicao2004}. Ela consome tempo, pois trata-se de ler, no mínimo, duas versões de texto de cada aluno. Ainda consome atenção, memória e raciocínio durante a leitura, a compreensão e a análise dos aspectos gramaticais e discursivos de cada produção textual. Sobretudo, engloba uma reflexão ampla e uma avaliação global do progresso de cada aluno por meio de comparações entre versões diferentes do texto.

A plataforma ``Redação Paraná'' foi desenhada para colaborar nas atividades de ensino de produção textual, isto é, contribuir para o trabalho do professor, mas ela também foi planejada para aprimorar o aprendizado da escrita e para o desenvolvimento da competência linguística dos alunos. Levando em consideração essas atribuições, pode-se afirmar que a ``Redação Paraná'' é relevante para essas atividades.

\textcite{conceicao2004} aponta que alguns professores de língua portuguesa podem ficar apenas na correção gramatical, a qual é superficial, e outros, ao tentarem analisar as questões discursivas do texto, acabam não sabendo como fazê-lo. Nesse contexto, o turno de correção da IA contribui ao efetuar uma revisão gramatical inicial, concomitantemente ao processo de escrita do texto pelo aluno.

Assim, ela realiza a correção indicativa e classificatória, indicando e classificando os problemas por meio de cores \cite{ruiz2015} no texto do aluno. Ainda proporciona uma correção resolutiva, já que a IA fornece a palavra ou expressão corrigida ao aluno. Como visto, no entanto, a IA não possui robustez para encontrar todos os problemas gramaticais e, assim, esses serão novamente aspectos que o professor, em seu turno, terá de verificar, o que é um aspecto limitante desta tecnologia.

No turno do professor, a existência de uma rubrica, isto é, de perguntas norteadoras e critérios de avaliação na ``Redação Paraná'', facilita a correção do professor em seu turno. Outro ponto positivo na plataforma está na possibilidade de fazer comentários ao final da correção, momento em que se realiza a correção textual-interativa. Nesta etapa, o professor mantém uma relação dialógica com o aluno e com seu texto, pois a correção textual-interativa favorece a reflexão discursiva do aluno sobre o seu texto, contribui para o seu aprimoramento linguístico na reescrita e ainda é uma oportunidade de ter um \textit{feedback} do seu professor sobre seu progresso como escritor/enunciador.

A partir desses apontamentos, pode-se afirmar que a plataforma ``Redação Paraná'' auxilia nas tarefas de ensino e de aprendizagem da escrita, colaborando nas atividades de correção do docente e nas atividades de escrita e reescrita do discente. No entanto, ainda apresenta aspectos limitantes, como a incapacidade de apontar todos os erros gramaticais do texto do aluno, bem como de caracterizá-los de maneira mais detalhada, a fim de permitir que o estudante compreenda a natureza dessas inadequações.

No que se refere aos limites deste artigo, é importante ressaltar a necessidade de se realizar pesquisas em campo sobre as práticas de escrita e de correção proporcionadas por estas plataformas em sala de aula. Apontaram-se aqui alguns pontos positivos e limitações, a partir de uma análise descritiva de suas funcionalidades, mas somente os usuários dessa tecnologia podem corroborar ou não esses aspectos.



\printbibliography\label{sec-bib}
% if the text is not in Portuguese, it might be necessary to use the code below instead to print the correct ABNT abbreviations [s.n.], [s.l.]
%\begin{portuguese}
%\printbibliography[title={Bibliography}]
%\end{portuguese}


%full list: conceptualization,datacuration,formalanalysis,funding,investigation,methodology,projadm,resources,software,supervision,validation,visualization,writing,review
\begin{contributors}[sec-contributors]
\authorcontribution{Agatha Suelin de Lima}[investigation,writing]
\authorcontribution{Maisa Morelatto}[investigation,writing]
\authorcontribution{Vanda Mari Trombetta}[conceptualization,investigation,methodology,projadm,visualization,writing,review]
\authorcontribution{Lovania Roehrig Teixeira}[conceptualization,investigation,visualization,writing,review]
\end{contributors}

\begin{dataavailability}
\txtdataavailability{nodata} % options: dataavailable, dataonly, databody, datanotav, nodata
\end{dataavailability}



\end{document}

